\section{ARQUITECTURA Y DISEÑO DE HARDWARE}

El hardware se construye alrededor de la placa OM13085 (LPCXpresso1769, revisión D2), que
incluye el microcontrolador, el cristal de 12\,MHz, el regulador de 3,3\,V, el LED RGB y la
sonda CMSIS-DAP para programación y depuración por USB. A la placa se conectan los
sensores, los mandos del usuario, la etapa de potencia del motor y el conversor serie. El
diagrama de bloques se incluye en el Anexo A.

\subsection{ETAPA DE ALIMENTACIÓN}

La placa de control (SSE.2) se alimenta desde el puerto USB de la sonda CMSIS-DAP; el
regulador de la propia placa genera los 3,3\,V del microcontrolador, que también alimentan
al potenciómetro. La tensión de referencia del ADC (VREFP) es la de 3,3\,V de la placa.

La etapa de potencia del motor y el buzzer trabajan con 5\,V.
\pendiente{indicar la fuente de 5\,V utilizada (USB, fuente externa, corriente máxima)
y la tensión de alimentación del LM35 y del sensor óptico}. Todas las masas son comunes.

\subsection{ETAPA DE CONTROL}

\subsubsection{DESCRIPCIÓN SSEP.2.1 (MICROCONTROLADOR)}

El SSEP.2.1 es el LPC1769 de NXP: núcleo ARM Cortex-M3, 512\,kB de memoria flash y 64\,kB
de SRAM. El firmware configura el PLL0 a partir del oscilador principal de 12\,MHz para
obtener un reloj de CPU (CCLK) de 100\,MHz, con un tiempo de acceso a flash de cinco
ciclos. Todos los periféricos usan PCLK = CCLK/4 = 25\,MHz.

La siguiente tabla resume los periféricos utilizados y su configuración.

\begin{longtable}{|p{2.4cm}|p{12.2cm}|}
\hline
\textbf{Periférico} & \textbf{Uso y configuración} \\ \hline
\endhead
TIMER0 & Base de tiempos de 1\,ms (prescaler PR = 24999). El match MR0 interrumpe cada
500 cuentas, sin reiniciar el contador, y cierra la ventana de medición de velocidad.
Interrupción con prioridad 1 (la más alta). \\ \hline
TIMER1 & Modo contador de flancos ascendentes en CAP1.0 (P1.18): cuenta los pulsos del
sensor óptico sin intervención de la CPU. \\ \hline
PWM1 & Contador a 1\,MHz (PR = 24) y MR0 = 1000: frecuencia de 1\,kHz con 1000 pasos de
resolución. El canal 2 (P2.1), en modo flanco simple, recibe el mando en MR2, de modo que
cada unidad por mil equivale a un paso. \\ \hline
ADC & Modo ráfaga sobre los canales AD0.0 (potenciómetro) y AD0.1 (LM35), 12 bits. Con
CLKDIV = 255 el reloj del ADC es de 97,7\,kHz: unas 1526 conversiones por segundo, 763
pares de muestras por segundo. \\ \hline
GPDMA & Canal 0 en modo periférico a memoria, disparado por el ADC, con una lista enlazada
circular de dos descriptores: AD0DR0 hacia la muestra del potenciómetro y AD0DR1 hacia la
del LM35. Interrumpe (prioridad 3) al completar cada par. \\ \hline
UART3 & 9600 baudios, 8N1, FIFO habilitadas. Recepción y transmisión por interrupción
(prioridad 2) con buffers circulares de 64 y 128 bytes. \\ \hline
GPIO & Entrada de la llave selectora (P2.6) y salidas del LED rojo (P0.22) y del buzzer
(P2.7). \\ \hline
\end{longtable}

La programación y depuración se realizan con la sonda CMSIS-DAP integrada en la placa.
La imagen Release compilada con MCUXpresso ocupa unos 10\,kB de flash y 360\,bytes de SRAM
estática, sin contar la pila, muy por debajo de la capacidad del microcontrolador.

\subsection{INTERFAZ DE USUARIO}

\textbf{Perilla (SSE.3.1).} Potenciómetro de 10\,k$\Omega$ conectado entre 3,3\,V y GND,
con el cursor en P0.23 (AD0.0). El ADC de 12 bits entrega de 0 a 4095 cuentas, que el
firmware convierte en un mando de 0 a 1000 por mil con redondeo al más cercano.

\textbf{Llave selectora (SSE.3.2).} Contacto entre P2.6 y GND, con la resistencia de
pull-up interna del pin habilitada. Llave abierta (nivel alto) corresponde a modo MANUAL y
llave cerrada (nivel bajo) a modo AUTOMÁTICO. Como el firmware la consulta una vez cada
500\,ms, mucho más que la duración del rebote de un contacto, no hace falta antirrebote
adicional.

\textbf{Alarma visual (SSE.3.3).} Se usa el LED rojo del LED RGB de la placa OM13085,
conectado a P0.22. Es de ánodo común, por lo que se enciende con nivel bajo.

\textbf{Alarma acústica (SSE.3.4).} Buzzer activo de 5\,V comandado por P2.7, activo por
alto. Como el pin entrega 3,3\,V y una corriente limitada, el buzzer se conecta mediante
una etapa de conmutación. \pendiente{describir la etapa del buzzer: transistor, resistencia
de base y diodo, si corresponde}.

\textbf{Interfaz en la notebook (SSE.3.5).} Una terminal serie (por ejemplo, PuTTY,
minicom o el monitor serie del IDE) configurada a 9600 baudios 8N1, en la que el usuario
escribe las consignas y lee la telemetría.

\subsection{ETAPA DE AISLACIÓN}

El diseño no incorpora una etapa de aislación galvánica: la placa de control, los sensores
y la etapa de potencia comparten la masa. Se consideró aceptable porque el motor es de
baja potencia y trabaja a 5\,V, y la etapa de potencia del motor ya separa al
microcontrolador de la corriente de carga. Como mejora, podría intercalarse un
optoacoplador entre P2.1 y la entrada del driver del motor, con alimentación separada para
la etapa de potencia.

\subsection{ETAPA DE POTENCIA}

El motor DC con reductora (SSE.5.1.1) se acciona mediante un driver (SSE.5.1) alimentado
con 5\,V, que recibe la señal PWM de 1\,kHz y 3,3\,V de P2.1.
\pendiente{indicar el modelo del driver o el circuito (transistor/MOSFET, diodo de rueda
libre) y la corriente del motor}.

Con el PWM al 100\,\% el motor alcanza aproximadamente 185\,RPM, según una medición
anterior del grupo; por eso la consigna máxima del modo automático es de 185\,RPM. Hasta
que el firmware configura el PWM, el pin P2.1 queda como entrada con pull-up (estado de
reset), por lo que la etapa de potencia debe tolerar ese estado sin arrancar el motor; el
firmware configura el PWM con mando cero antes de cualquier otra salida.

\subsection{ETAPA DE ADQUISICIÓN}

\textbf{Sensor de temperatura (SSE.6.1).} LM35 conectado a P0.24 (AD0.1). Su salida es de
10\,mV/\textdegree C. Con VREFP = 3,3\,V el ADC de 12 bits resuelve 0,806\,mV por cuenta,
es decir unos 0,08\,\textdegree C por cuenta. La conversión se hace en aritmética entera,
en centésimas de grado:
\[
T_{[0{,}01\,^\circ\mathrm{C}]} = \mathrm{redondeo}\!\left(\frac{N \cdot 3300 \cdot 10}{4096}\right)
\]
donde $N$ es el resultado del ADC. El firmware promedia las últimas cuatro lecturas (una
cada 500\,ms) para que el ruido del ADC no haga oscilar la alarma. El anteproyecto preveía
una referencia externa de 2,54\,V; finalmente se usa la de 3,3\,V de la placa, lo que reduce
levemente la resolución pero simplifica el hardware. El valor de VREFP se configura al
compilar (\texttt{ADC\_VREF\_MV}) y puede ajustarse al valor medido.

\textbf{Sensor óptico de ranura (SSE.6.2).} Módulo XD-51, que reemplaza al HC-89 previsto
en el anteproyecto, montado sobre un disco de 60 ranuras solidario al eje. Su salida se
conecta a P1.18 (CAP1.0), configurado con pull-up interno porque la salida de estos
módulos suele ser un comparador a colector abierto. TIMER1 cuenta un flanco ascendente por
ranura, es decir 60 pulsos por vuelta. La velocidad se obtiene como:
\[
\mathrm{RPM} = \frac{\text{pulsos en la ventana}}{60\ \text{pulsos/vuelta}} \cdot
\frac{60\,000\ \mathrm{ms/min}}{500\ \mathrm{ms}} = 2 \cdot \text{pulsos}
\]
La resolución es, por lo tanto, de 2\,RPM. A la velocidad máxima de 185\,RPM el sensor
genera 185 pulsos por segundo, muy por debajo del límite del contador (PCLK/4).

\subsection{ETAPA DE COMUNICACIÓN}

La comunicación con la notebook (SSE.7) usa UART3 del microcontrolador, en P0.0 (TXD3) y
P0.1 (RXD3), con niveles de 3,3\,V, y un conversor USB--UART (SSE.7.1) que la notebook ve
como un puerto serie virtual. \pendiente{indicar el modelo del conversor (CP2102, CH340,
FT232, etc.) y verificar que trabaje con niveles de 3,3\,V}. El pin de recepción tiene
pull-up interno para que la línea quede en reposo aunque el conversor esté desconectado.

A 9600 baudios, la biblioteca calcula DLL = 93 y un divisor fraccional 3/4, que dan una
velocidad real de 9600,6 baudios (error menor a 0,01\,\%).
